%%
%% Front matter: rights, conference, title, authors, abstract, keywords.
%% Included by main.tex before \maketitle.

%%
%% Rights management information.
\copyrightyear{2026}
\copyrightclause{Copyright for this paper by its authors.
  Use permitted under Creative Commons License Attribution 4.0
  International (CC BY 4.0).}

%%
%% Conference information.
\conference{CLEF 2026 -- Conference and Labs of the Evaluation Forum,
  September 21--24, 2026, Jena, Germany}

%%
%% Title.
\title{Partial Fine-Tuning of BioCLIP 2.5 with Taxonomic Auxiliary Heads for Multi-Label Plant Species Identification in Vegetation Plots: ANU at PlantCLEF 2026}

%%
%% Authors.
%% Per the canonical CEUR-ART template, each author block supports an
%% optional key-value list with \texttt{orcid=}, \texttt{email=},
%% \texttt{url=}. Only \texttt{email=} is supplied here; add the other
%% keys when convenient.
\author[1]{Arjun Raj}[%
  email=u7526852@anu.edu.au,
]
\fnmark[1]

\address[1]{The Australian National University, Canberra, ACT 2601, Australia}

\author[1]{Manindra de Mel}[%
  email=u7543337@anu.edu.au,
]
\cormark[1]
\fnmark[1]

\author[1]{Razeen Wasif}[%
  email=u7619822@anu.edu.au,
]
\fnmark[1]

\author[1]{William Brake}[%
  email=u7480294@anu.edu.au,
]
\fnmark[1]

\cortext[1]{Corresponding author.}
\fntext[1]{These authors contributed equally.}

%%
%% Abstract.
\begin{abstract}
This paper describes the system submitted by the Australian National University team to the PlantCLEF 2026 challenge, which tasks participants with predicting all plant species present in high-resolution vegetation quadrat images. The core difficulty lies in the domain shift between the training data, comprising individually photographed single-plant specimens, and the test data, consisting of dense multi-species vegetation plots. A secondary challenge is the long-tail distribution of the training data, where common species have thousands of images while rare species have fewer than ten. Our final submission centres on a partial fine-tune of BioCLIP 2.5 (ViT-H/14) in which only the last four transformer blocks together with the final layer norm and projection are unfrozen, while the lower twenty-eight blocks are kept frozen to preserve the Tree-of-Life prior. The encoder is trained on the enlarged PlantCLEF 2024 single-plant manifest with auxiliary cross-entropy heads at the genus and family levels, supplying a structural regulariser without adding inference-time cost. Inference decomposes each quadrat into a non-overlapping $4{\times}4$ grid of sixteen tiles, forwards each tile through the encoder at both $224$ and $336$ pixels, and aggregates the tile distributions by per-species softmax mean across the sixteen tiles and across the two resolutions; class-prior logit adjustment ($\tau{=}0.25$) is then applied and an adaptive probability threshold ($T{=}0.03$, $k_{\min}{=}2$, $k_{\max}{=}10$) selects the emitted species set. Beyond this anchor we report three architecture-orthogonal pivots that we submitted but that did not improve over the visual posterior: a triple-backbone classifier retrained on frozen BioCLIP 2.5, DINOv3, and ConvNeXt-V2-L features; a genus and family co-occurrence rerank; and a seasonal phenology prior derived from GBIF month-of-observation histograms. Across all four submitted configurations, our best public Macro F1 is \textbf{0.418265} and our best private Macro F1 is \textbf{0.40283}.
\end{abstract}

%%
%% Keywords.
\begin{keywords}
  plant identification \sep
  multi-label classification \sep
  partial fine-tuning \sep
  BioCLIP \sep
  taxonomic auxiliary loss \sep
  vegetation plot analysis \sep
  PlantCLEF 2026
\end{keywords}
